\chapter{Conclusiones}
\label{cap:Conclusiones}

\section{Objetivos alcanzados}

Se ha alcanzado el objetivo principal del trabajo: el diseño e implementación de un prototipo funcional de codificación clínica automática en el estándar CIE-10-ES. El sistema resultante integra un clasificador multilabel basado en \emph{transformers} con una arquitectura completa de software que permite a un codificador clínico enviar casos clínicos en español y recibir propuestas de codificación con una confianza relativa asociada a cada código. La propuesta requiere siempre la revisión completa de un codificador clínico, y los resultados respaldan la viabilidad técnica del enfoque, no su utilidad clínica, que no se ha evaluado con usuarios.

A continuación se revisa el grado de cumplimiento de cada uno de los subobjetivos enumerados en el capítulo~\ref{cap:Objetivo}:

\begin{enumerate}
 \item \textbf{Adquisición y construcción de las fuentes de datos.} Completado. Se descargó el corpus CodiEsp en sus tres particiones (diagnósticos, procedimientos y variante extendida con anotaciones a nivel de \emph{span}), y se desarrollaron los scripts de \emph{scraping} sobre el portal eCIE-maps del Ministerio de Sanidad para obtener el catálogo completo de códigos CIE-10-ES, cuyo formato de descarga pública estructurada no existe. El resultado son los ficheros CSV que alimentan todas las fases posteriores del proyecto. Se implementó el \emph{pipeline} de carga, limpieza y transformación de las anotaciones al formato multilabel requerido por el modelo, incluyendo la conversión de las anotaciones de nivel \emph{span} al nivel de documento empleado en el entrenamiento.


 \item \textbf{Análisis estadístico de las fuentes de datos.} Completado. El análisis de la distribución de longitudes, la frecuencia de aparición de los códigos y el desequilibrio de clases fundamentó las principales decisiones de diseño del modelo: el uso de ponderación de clases en la función de pérdida y la exploración de umbrales de decisión por código, finalmente descartados en favor de un umbral global (véase el subobjetivo siguiente).

 \item \textbf{Módulo clasificador multilabel.} Completado. Tras la comparativa de diez modelos candidatos, cinco evaluados (Anexo~\ref{cap:AnexoA}), se seleccionó \texttt{IIC/RigoBERTa-Clinical} como modelo base y se ajustó mediante \emph{fine-tuning} sobre la partición de diagnósticos de CodiEsp, que es el alcance definido en el capítulo de objetivos; la codificación de procedimientos queda fuera del trabajo. El modelo final, entrenado sobre 500 documentos con 1.767 códigos objetivo, es el clasificador del sistema (umbral global de 0,2 con el clasificador solo y umbral propio sobre la puntuación combinada en el modo fusionado; los umbrales por clase se descartaron por sobreajustar la validación). Sobre el conjunto de prueba alcanza un \textbf{F1-micro de 0,495} y un \textbf{F1-macro de 0,108}, calculados sobre los pares anotados cuyo código pertenece al vocabulario del modelo; sobre todos los pares anotados el F1-micro es 0,444, porque el modelo no puede acertar códigos que no vio en el entrenamiento. La brecha entre F1-micro y F1-macro refleja el fuerte desequilibrio de clases del corpus, donde muchos códigos cuentan con muy pocos ejemplos de entrenamiento.

La calidad del \emph{ranking} de códigos, medida como \textbf{MAP por documento}, es de \textbf{0,454} sobre el \emph{gold} completo con el clasificador solo, directamente comparable con el \emph{benchmark} CodiEsp-D 2020~\cite{miranda2020}. Con el clasificador solo queda en el penúltimo lugar de los diez sistemas comparados, solo por encima de IMS, lo que es coherente con un clasificador plano entrenado con 500 documentos. La cifra de 0,555 corresponde al MAP restringido al vocabulario que el modelo puede predecir y no es comparable con el \emph{benchmark}, cuyos sistemas podían proponer cualquier código.

El \textbf{modo fusionado}, que es el que sirve el sistema por defecto, combina el clasificador con el diccionario de patrones clínicos en un único \emph{ranking} y alcanza \textbf{MAP 0,554 y la segunda posición}, por delante del resto de sistemas, incluidos los diccionarios construidos a mano y los basados en \emph{transformers}, y solo por detrás de IXA-AAA (MAP~0,593). Los resultados sugieren que las dos fuentes de evidencia fallan en sitios distintos: el diccionario acierta el bloque de tres caracteres y no sabe ordenar dentro de él, mientras que el clasificador ordena mejor dentro del bloque pero elige peor el bloque.

 \item \textbf{Backend de gestión y trazabilidad.} Completado. Se desarrolló un backend en Elixir con Phoenix que orquesta el flujo de análisis, gestiona la autenticación con confirmación por correo electrónico y aplica una arquitectura CQRS/\emph{Event Sourcing} sobre PostgreSQL, de modo que cada acción clínica queda registrada como un evento inmutable con su marca temporal y su autoría.

 \item \textbf{Servicio de inferencia (API REST).} Completado. El motor de IA expone mediante FastAPI los \emph{endpoints} de predicción, explicación, resumen y estado (Anexo~\ref{cap:AnexoM}), con una imagen Docker optimizada para PyTorch y soporte para CPU y GPU. El servicio arranca incluso sin modelo entrenado, con un código 503 informativo, lo que facilita el desarrollo y las pruebas.

 \item \textbf{Módulo de generación de resúmenes.} Completado. Se integró un modelo de lenguaje generativo ligero (cuantizado en formato GGUF, ejecutable en CPU) que produce un resumen o una paráfrasis del informe, con modelo, modo y \emph{prompts} configurables desde el \emph{backoffice} y emitido en tiempo real mediante \emph{streaming}. Por tratarse de una ayuda opcional y por su coste computacional, permanece desactivado por defecto en el despliegue y se habilita bajo demanda. Su calidad se validó de forma cualitativa durante el desarrollo; una evaluación cuantitativa sistemática de los resúmenes generados (por ejemplo, con métricas de fidelidad factual y de cobertura de la información clínica) queda pendiente como trabajo futuro.

 \item \textbf{Interfaz web de demostración.} Completada. La interfaz web, desarrollada en Next.js 16 con React 19, permite enviar informes clínicos y visualizar en tiempo real las tarjetas de análisis generadas por el sistema. El codificador puede validar o rechazar individualmente cada código predicho y consultar el historial de análisis anteriores desde el panel lateral.
\end{enumerate}

Más allá de los subobjetivos enumerados, el desarrollo ha incluido componentes complementarios: el \emph{backoffice} de administración, la internacionalización en cinco idiomas, y una infraestructura de despliegue con Docker Compose, Traefik y un túnel de Cloudflare.

\section{Verificación de los requisitos}

La Tabla~\ref{tab:verificacion-req} resume el grado de cumplimiento de los requisitos definidos en la sección~\ref{sec:requisitos} del capítulo~\ref{cap:Desarrollo}.

{\footnotesize
\begin{longtable}{llP{10.4cm}}
\caption{Verificación de los requisitos funcionales (RF) y no funcionales (RNF).}\label{tab:verificacion-req}\\
\hline
\textbf{Req.} & \textbf{Estado} & \textbf{Evidencia} \\
\hline
\endfirsthead
\hline
\textbf{Req.} & \textbf{Estado} & \textbf{Evidencia} \\
\hline
\endhead
RF1 & Cumplido & El motor acepta texto clínico libre sin marcado previo (prueba de \texttt{POST /predict} en \texttt{ai\_engine/\allowbreak{}tests/\allowbreak{}test\_\allowbreak{}api.py}). \\
RF2 & Cumplido & \texttt{POST /predict} devuelve los códigos de diagnóstico predichos (\texttt{test\_\allowbreak{}api.py} y prueba de extremo a extremo de envío de un informe en \texttt{e2e/\allowbreak{}test\_\allowbreak{}flujos.py}). \\
RF3 & Cumplido & Cada código incluye su probabilidad y una confianza relativa normalizada, no calibrada (sección~\ref{sec:confianza-predicciones}); la presencia del campo la comprueba \texttt{test\_\allowbreak{}api.py}. \\
RF4 & No cumplido & El servicio procesa un informe por petición y no existe un \emph{endpoint} que reciba una lista de textos y devuelva los resultados en el mismo orden. Un cliente puede procesar un lote con llamadas independientes a la API (también concurrentes), pero eso no cumple el requisito tal como está formulado. \\
RF5 & Cumplido & API REST documentada con OpenAPI, con \emph{endpoint} de inferencia y de estado (Anexo~\ref{cap:AnexoM}). La especificación la genera FastAPI y no tiene una prueba automática propia; el \emph{endpoint} de estado sí lo cubre \texttt{test\_\allowbreak{}api.py}. \\
RF6 & Cumplido & \texttt{POST /auth/register} da de alta sin conceder acceso y envía un correo con el enlace de \texttt{GET /auth/confirm/:token}; el recorrido completo lo cubre la prueba de extremo a extremo \texttt{test\_confirmar\_el\_correo\_da\_acceso}. \\
RF7 & Cumplido & El canal WebSocket emite \texttt{analysis\_card\_received} con los códigos y, después, \texttt{triggers\_received} con los términos, de modo que la atribución no retrasa la aparición de los códigos (Anexo~\ref{cap:AnexoC}). \\
RF8 & Cumplido & Los eventos \texttt{CodeValidated} y \texttt{CodeRejected} registran la decisión y, en el rechazo, su motivo; la interfaz exige el motivo antes de habilitar la confirmación. \\
RF9 & Cumplido & \texttt{POST /summarize/stream} devuelve el resumen en \emph{streaming} con un modelo cuantizado ejecutable en CPU (prueba del \emph{endpoint} sin modelo en \texttt{test\_\allowbreak{}api.py}; la calidad de los resúmenes no se ha medido). \\
RF10 & Cumplido & Interfaz en Next.js con envío de informes, seguimiento del análisis, términos por código e histórico de análisis con papelera (pruebas de componentes \texttt{Chat\allowbreak{}Interface.\allowbreak{}test.\allowbreak{}tsx} y \texttt{Conversation\allowbreak{}Sidebar.\allowbreak{}test.\allowbreak{}tsx} en \texttt{frontend/\allowbreak{}src/\allowbreak{}test/\allowbreak{}components}). \\
\hline
RNF1 & Cumplido & MAP 0,554 en el modo fusionado servido por defecto (segundo de la comparativa) y 0,454 con el clasificador solo (penúltimo de diez); Tabla~\ref{tab:codiesp-comparativa}. \\
RNF2 & Cumplido & Por informe, en CPU (Ryzen 9 7900X) y en GPU (RTX 4080): predecir 0,331\,s y 0,017\,s; explicar 16,2\,s y 0,51\,s; ciclo completo 16,5\,s y 0,53\,s (Tabla~\ref{tab:latencia}). El requisito acota la predicción, que es lo que el codificador clínico espera para ver los códigos; la explicación se sirve aparte y llega después. \\
RNF3 & Cumplido & Servicios desacoplados en contenedores con interfaces explícitas; el modelo base se sustituye sin tocar el resto. \\
RNF4 & Cumplido & Cada predicción emite un evento \texttt{AIPredictionReceived} con el identificador del informe, el motor empleado y la versión del modelo, entendida como el nombre del \emph{encoder} y el hash del fichero de pesos. \\
RNF5 & Cumplido & Exportación y borrado de cuenta, cifrado en tránsito y minimización de datos (sección~\ref{sec:rgpd}). \\
\hline
\end{longtable}
}

Un requisito no se ha satisfecho, y se reconoce de forma explícita: la \textbf{inferencia por lotes} (RF4) quedó fuera del alcance final por no ser necesaria para el prototipo de demostración, y queda como mejora pendiente. El \textbf{rendimiento} (RNF1) se cumple: el modo fusionado servido por defecto es segundo de la comparativa (MAP 0,554).

\section{Justificación de competencias adquiridas}

En el TFG se han aplicado las competencias correspondientes a la Tecnología Específica de \emph{Computación}:

\begin{description}
 \item[CO4:] \emph{Capacidad para conocer y desarrollar técnicas de aprendizaje computacional y diseñar e implementar aplicaciones y sistemas que las utilicen, incluyendo las dedicadas a extracción automática de información y conocimiento a partir de grandes volúmenes de datos.} Esta competencia es el núcleo del trabajo: se ha diseñado, implementado y evaluado un clasificador multilabel de texto clínico basado en \emph{transformers}, realizando el ciclo completo de aprendizaje supervisado por transferencia (\emph{fine-tuning}) sobre el corpus CodiEsp. El proceso ha requerido la aplicación de técnicas de aprendizaje profundo: función de pérdida Binary Cross-Entropy con ponderación de clases para combatir el desequilibrio y un término de ordenación (ZLPR), planificación de la tasa de aprendizaje con \emph{cosine schedule} y \emph{plateau scheduler}, descongelación progresiva de capas del \emph{encoder}, y selección de un umbral de decisión global sobre el conjunto de validación (los umbrales por clase se probaron y se descartaron). El sistema resultante extrae automáticamente códigos CIE-10-ES a partir de informes clínicos en texto libre, y ha sido evaluado con métricas estándar de clasificación multilabel (F1-micro y F1-macro) y con el MAP por documento, la métrica del \emph{benchmark} CodiEsp-D 2020 de referencia en español.
\end{description}

\section{Trabajos derivados y futuros}

Seis líneas de trabajo concretas surgen del estado actual del prototipo.

\textbf{Ampliación del corpus de entrenamiento}. El clasificador se ha entrenado sobre 500 documentos de CodiEsp, lo que limita la cobertura de los códigos de enfermedades raras o procedimientos poco frecuentes, como refleja la diferencia entre el F1-micro de 0,495 (0,444 sobre todos los pares anotados) y el F1-macro de 0,108 sobre el conjunto de prueba (umbral global de 0,2). A esta limitación se añade una restricción estructural del espacio de salida, inherente a los datos y no al modelo: el clasificador solo puede predecir los 1.767 códigos de la partición de entrenamiento de CodiEsp (el corpus completo tiene 2.557), en torno al 2,5\,\% de los 69.823 del catálogo CIE-10-ES completo. El resto de códigos no aparece en el corpus público: un clasificador supervisado plano no puede asignar un código ausente de su espacio de salida. Ampliar ese espacio exige ejemplos anotados o bien enfoques \emph{zero-shot} apoyados en las descripciones del catálogo, línea que este trabajo solo explora de forma parcial mediante el preentrenamiento débil. Una vía para ampliar esa cobertura es el ciclo de aprendizaje continuo descrito en la sección~\ref{sec:continuous-learning}: si los codificadores validan o corrigen las predicciones en uso real, cada análisis podría aportar ejemplos etiquetados de nuevos códigos y ampliar el vocabulario del modelo con el uso. El F1-micro promedia ponderando por frecuencia de código, por lo que los códigos comunes dominan el resultado; el F1-macro da el mismo peso a cada código y penaliza directamente los raros con pocos ejemplos de entrenamiento. La incorporación de historias clínicas reales anonimizadas, en colaboración con centros hospitalarios, permitiría aumentar la representación de esos códigos y mejorar la utilidad clínica del sistema precisamente en los casos más atípicos. Esta colaboración parece una de las vías con mayor margen de mejora.

\textbf{Los códigos residuales exigen razonar sobre lo que el informe no dice}. El análisis de error del conjunto de prueba (Tabla~\ref{tab:analisis-error}) deja al descubierto una limitación que probablemente no se resuelva solo con más datos. Los códigos que el modelo más sobrepredice son también los que más deja escapar: R52 «dolor, no especificado» encabeza los falsos positivos y aparece también entre los falsos negativos, junto a R69 «enfermedad, no especificada»; entre los falsos positivos más frecuentes figuran además N28.9 «trastorno del riñón y del uréter, no especificado» y R53.1 «astenia». Todos ellos son residuales. Son las categorías que el manual de codificación reserva para cuando el informe menciona un hallazgo sin concretarlo, de modo que asignarlas correctamente depende de comprobar que el documento no aporta el detalle que llevaría a un código más específico. Un clasificador plano que puntúa cada etiqueta a partir de lo que sí aparece en el texto no tiene forma de representar esa ausencia, y por eso los confunde en ambos sentidos. La vía que sí la aborda es reformular la salida como una decisión jerárquica sobre el árbol CIE-10, en la que el modelo elija primero el bloque y solo después descienda al código, con el residual como resultado explícito de no encontrar evidencia para ninguna rama más fina. Queda fuera del alcance de este trabajo porque cambia la arquitectura de la cabeza clasificadora y el esquema de evaluación.

\textbf{Migración a infraestructura gestionada en la nube}. El sistema se sirve desde una estación de trabajo propia con Docker Compose y Traefik, publicada mediante un túnel de Cloudflare, sin proveedor en la nube. Si el volumen de usuarios o la latencia de inferencia lo exigieran, podría trasladarse a un entorno con GPU (nube gestionada, Kubernetes propio o instancia dedicada); al estar dockerizado, es previsible que requiera pocos cambios, pero no se ha probado. El Anexo~\ref{cap:AnexoI} recoge la infraestructura actual y la estimación de costes de esa alternativa.

\textbf{Explorar la combinación de varias ejecuciones}. Al medir la variabilidad estocástica del entrenamiento se observó que promediar las salidas de tres ejecuciones disponibles reducía la varianza y mejoraba el MAP sobre el conjunto de prueba (sección~\ref{sec:anexoF-map}). Se trata de una observación incidental sobre ejecuciones que perseguían otros objetivos, no de un sistema diseñado ni evaluado como tal, por lo que no forma parte de los resultados de este trabajo. Diseñar y evaluar una combinación de modelos, midiendo su coste de inferencia frente a la ganancia, es una línea de trabajo abierta. El obstáculo para llevarla a producción es precisamente ese coste: el sistema debe poder servirse sin acelerador, y en modo CPU multiplicar las pasadas del \emph{encoder} por informe no resulta viable. La vía que sortearía esa restricción es la \textbf{destilación de conocimiento}~\cite{hinton2015distilling}: entrenar un único clasificador contra las probabilidades producidas por la combinación, en lugar de contra las etiquetas binarias del corpus, conservando así el coste de inferencia de un solo modelo. Una prueba preliminar (sección~\ref{sec:anexoF-map}) muestra que esa vía no es inmediata en este régimen de datos: con 500 documentos de entrenamiento el profesor memoriza el corpus, sus probabilidades sobre él coinciden con las etiquetas binarias y el alumno no hereda nada del conjunto. Lo que queda abierto es destilar sobre un corpus clínico no etiquetado y ajeno al entrenamiento, donde el profesor sí produciría probabilidades informativas; requiere textos adicionales de los que este trabajo no dispone.

\textbf{Mejora de la explicabilidad}. El sistema atribuye la importancia de cada palabra por perturbación: sustituye la palabra por la máscara del \emph{tokenizador} y mide cuánto cae el \emph{logit} del código; un filtro por gradiente limita el cálculo a las 32 palabras más prometedoras, porque el gradiente por sí solo producía atribuciones ruidosas en una arquitectura de 24 capas (sección~\ref{sec:explicabilidad}). El método desplegado es fiel por construcción, porque mide el efecto real de quitar cada palabra, pero su coste crece con la longitud del informe y no cuantifica la incertidumbre de la atribución. Las dos extensiones naturales son \emph{SmoothGrad}~\cite{Smilkov17}, que promedia varias perturbaciones con ruido para estimar intervalos sobre la importancia, y \emph{SHAP KernelExplainer}~\cite{Lundberg17}, que reparte la importancia entre palabras. Ambas añaden coste por consulta, de modo que la pregunta que queda abierta es empírica: si la ganancia en calidad de la explicación justifica ese coste, medido sobre un conjunto de informes anotados con los términos clínicamente relevantes. Falta calibrar la confianza mostrada (sección~\ref{sec:confianza-predicciones}), por ejemplo con escalado de temperatura.

\textbf{Despliegue en entorno hospitalario real con captura de retroalimentación de los codificadores}. El prototipo incluye ya los mecanismos de validación y rechazo de códigos por parte del codificador, lo que sienta las bases para un ciclo de retroalimentación activa: las correcciones realizadas por los codificadores clínicos pueden usarse para reetiquetar documentos y realimentar el entrenamiento del modelo de forma iterativa. La puesta en marcha de este ciclo en un entorno real permitiría evaluar la utilidad del sistema en condiciones de uso auténticas y obtener datos de entrenamiento de mayor calidad que los disponibles en el corpus público.

Esa vía está construida pero sin uso: el codificador puede validar o rechazar cada código y marcar qué términos del informe justifican la predicción, y todo ello se persiste, pero la base de datos del prototipo registra cero códigos validados y cero términos verificados porque el sistema no se ha usado en un entorno clínico. Reentrenar el \emph{encoder} con esas correcciones depende por tanto de acumular uso real y no de escribir más código.

Esa misma restricción es la que llevó a \textbf{descartar} la otra línea que se llegó a plantear: el ajuste fino de un modelo generativo, del tipo Gemma, para predecir los códigos directamente en lugar de clasificarlos. No se ensayó; la decisión se apoya en dos resultados propios: la aumentación por retrotraducción, que cuadruplicó el corpus, empeoró el resultado, y la destilación falló porque con 500 documentos el modelo profesor memoriza el corpus. Con ese cuello de botella en los datos, es previsible que un modelo con más parámetros tenga poco recorrido sobre el mismo material, además de multiplicar el coste de inferencia en un sistema que debe poder servirse también sin acelerador. Solo volvería a tener sentido con un corpus anotado sustancialmente mayor.

\section{Viabilidad económica}

El sistema es un prototipo académico y no se ha realizado ningún estudio de mercado ni análisis regulatorio que respalde su comercialización. El despliegue actual no tiene coste directo, porque la estación de trabajo ya existía y se alimenta con autoconsumo solar; a precio de red, su consumo eléctrico sería de unos 330\,€/año (Anexo~\ref{cap:AnexoI}, que estima también el coste de la alternativa en la nube). Una explotación futura dependería sobre todo del soporte, del cumplimiento normativo y de la responsabilidad sobre un sistema que interviene en un proceso clínico, costes que no se han dimensionado. Tampoco se ha medido cuántas peticiones concurrentes soporta el despliegue actual, que no ha tenido usuarios reales. Por último, el ciclo de retroalimentación de la sección~\ref{sec:continuous-learning} podría aportar datos etiquetados por profesionales, pero hoy no hay ninguna validación registrada, de modo que su efecto sobre el modelo es una hipótesis.

\section{Valoración personal}

El proyecto ha supuesto un recorrido completo por el ciclo de vida de un sistema de \emph{software} con componente de investigación aplicada: desde la adquisición y el análisis de los datos hasta el despliegue, pasando por la selección y el entrenamiento del modelo, el diseño de la arquitectura y el desarrollo de todos los componentes del sistema.

La parte más exigente y con mayor desviación respecto a la planificación original ha sido el ciclo de experimentación del modelo (capítulo~\ref{cap:Planificacion} y Anexo~\ref{cap:AnexoF}). Esa experiencia ha permitido desarrollar una intuición práctica sobre el impacto de decisiones como la tasa de aprendizaje, la ponderación de clases o el umbral de decisión, que difícilmente se adquiere sin experimentación directa y observación de las curvas de aprendizaje.

Las herramientas de apoyo empleadas, por uso, fueron las siguientes:
\begin{itemize}
 \item \textbf{Código}: GitHub Copilot, como asistente de autocompletado, para agilizar bloques repetitivos y esqueletos de código. Claude Code (Anthropic) se empleó también en cambios puntuales de código y sus pruebas. El código se revisó mediante lectura y mediante el CI.
 \item \textbf{Memoria, revisión}: Claude Code (Anthropic), mediante una \emph{skill} de revisión de uso local: lectura crítica del texto, detección de incoherencias numéricas, terminológicas y de formato, contraste de las cifras con los ficheros de resultados, y propuestas puntuales de corrección de formato y de redacción.
 \item \textbf{Memoria, texto}: \texttt{aspell} para la corrección ortográfica y el complemento Grammarly para la corrección gramatical y de estilo.
\end{itemize}
Todo se revisó antes de incorporarlo y cada hallazgo se verificó contra los datos antes de aplicar ningún cambio; el contenido, las mediciones y las conclusiones son del autor.

El proyecto también ha puesto de manifiesto el valor de mantener los componentes del sistema con interfaces bien definidas. El uso del \emph{mock} del motor de IA para el desarrollo del backend y el frontend, con independencia del estado del modelo entrenado, ha evitado bloqueos entre las distintas líneas de trabajo y ha permitido avanzar en paralelo de forma efectiva, reduciendo el impacto de las iteraciones del ciclo de entrenamiento sobre el resto del sistema.

\enlargethispage{\baselineskip}
En términos cuantitativos, el proyecto se ha desarrollado a lo largo de aproximadamente 376 horas de trabajo efectivo, repartidas en 72 puntos de historia entre las seis épicas descritas en el capítulo~\ref{cap:Planificacion} y otros 22 dedicados a la redacción de la presente memoria, 94 puntos en total. El historial de desarrollo queda registrado en el repositorio público del proyecto~\cite{cie10repo2026}. Los primeros experimentos de entrenamiento se iniciaron en noviembre de 2024 en la rama \texttt{master}. En enero de 2026, aprovechando los scripts, análisis y conclusiones ya obtenidos, se reescribió el historial Git para desarrollar el sistema completo desde una base limpia. La rama \texttt{main} resultante contiene más de 360 \emph{commits}.
